\documentclass[11pt]{article}
\usepackage[margin=1in]{geometry}
\usepackage{amsmath, amssymb, amsthm}
\usepackage{graphicx}
\usepackage{hyperref}
\usepackage{cite}
\usepackage{fancyhdr}
\usepackage{setspace}

% Header setup
\pagestyle{fancy}
\fancyhf{}
\rhead{Reading Summary}
\lhead{Liu Zhonglin}
\cfoot{\thepage}

\title{Reading Summary: Simulation study in Probabilistic Boolean Network models for genetic regulatory networks}
\author{Liu Zhonglin \\ MATH4999 Mathematics project}
\date{\today}

\begin{document}
	
	\maketitle
	\thispagestyle{empty}
	
	
	
	\section{Executive Summary}
	\textit{[2-3 sentences providing a high-level overview of what the paper is about and its main contribution]}
	
	\section{Research Problem and Motivation}
	\subsection{Problem Statement}
	\textit{[What specific problem does this paper address?]}
	
	\subsection{Motivation and Significance}
	\textit{[Why is this problem important? What gap does it fill?]}
	
	\subsection{Research Questions/Objectives}
	\begin{itemize}
		\item \textit{[Primary research question]}
		\item \textit{[Secondary questions, if any]}
	\end{itemize}
	
	\section{Background and Related Work}
	\subsection{Theoretical Foundation}
	\textit{[Key concepts, theories, or frameworks the paper builds upon]}
	
	\subsection{Previous Work}
	\textit{[How does this work relate to existing literature? What are the key references?]}
	
	\subsection{Gaps in Existing Research}
	\textit{[What limitations or gaps in previous work does this paper address?]}
	
	\section{Methodology}
	\subsection{Approach Overview}
	\textit{[High-level description of the methodology]}
	
	\subsection{Key Methods/Techniques}
	\begin{itemize}
		\item \textit{[Method 1: brief description]}
		\item \textit{[Method 2: brief description]}
		\item \textit{[Additional methods as needed]}
	\end{itemize}
	
	\subsection{Data and Experimental Setup}
	\textit{[What data was used? How were experiments designed?]}
	
	\subsection{Evaluation Metrics}
	\textit{[How was success measured?]}
	
	\section{Key Contributions and Results}
	\subsection{Main Contributions}
	\begin{enumerate}
		\item \textit{[First major contribution]}
		\item \textit{[Second major contribution]}
		\item \textit{[Additional contributions]}
	\end{enumerate}
	
	\subsection{Experimental Results}
	\textit{[Summary of key findings and results]}
	
	\subsection{Performance Analysis}
	\textit{[Quantitative results, comparisons with baselines, etc.]}
	
	\section{Technical Details}
	\subsection{Algorithms/Models}
	\textit{[Key technical innovations, mathematical formulations]}
	
	\subsection{Implementation Details}
	\textit{[Important implementation aspects, if relevant]}
	
	\subsection{Complexity Analysis}
	\textit{[Computational complexity, scalability considerations]}
	
	\section{Strengths and Limitations}
	\subsection{Strengths}
	\begin{itemize}
		\item \textit{[Methodological strengths]}
		\item \textit{[Technical innovations]}
		\item \textit{[Experimental rigor]}
	\end{itemize}
	
	\subsection{Limitations}
	\begin{itemize}
		\item \textit{[Methodological limitations]}
		\item \textit{[Scope limitations]}
		\item \textit{[Experimental limitations]}
	\end{itemize}
	
	\section{Impact and Applications}
	\subsection{Practical Applications}
	\textit{[How can these results be applied in practice?]}
	
	\subsection{Broader Impact}
	\textit{[What is the significance for the field?]}
	
	\subsection{Future Research Directions}
	\textit{[What follow-up work does this enable or suggest?]}
	
	\section{Critical Analysis}
	\subsection{Validity of Claims}
	\textit{[Are the claims well-supported by the evidence?]}
	
	\subsection{Methodological Soundness}
	\textit{[Are the methods appropriate and well-executed?]}
	
	\subsection{Significance Assessment}
	\textit{[How significant is this contribution to the field?]}
	
	\section{Key Takeaways}
	\begin{itemize}
		\item \textit{[Most important lesson/insight]}
		\item \textit{[Key technical concept learned]}
		\item \textit{[Methodological insight]}
		\item \textit{[Relevance to your work/interests]}
	\end{itemize}
	
	\section{Questions and Discussion Points}
	\begin{itemize}
		\item \textit{[Unclear aspects that need clarification]}
		\item \textit{[Potential criticisms or concerns]}
		\item \textit{[Ideas for extending the work]}
	\end{itemize}
	
	\section{Relevance to Current Work}
	\textit{[How does this paper relate to your current research/project? What can you apply or build upon?]}
	
	\section{Additional Notes}
	\textit{[Any other observations, connections to other papers, or miscellaneous thoughts]}
	
	% Optional: Include key figures or equations from the paper
	% \section{Key Figures/Equations}
	% \textit{[Reproduce or describe important visual elements]}
	
	\bibliographystyle{plain}
	\nocite{*}
	
	\bibliography{ref}
	% Note: Create a references.bib file with relevant citations
	
\end{document}